\documentclass{article}
\usepackage[a4paper]{geometry}
\usepackage{fancyhdr}
\pagestyle{fancy}
\lhead{Artenerfassung}
\rhead{Mai 2026}

\begin{document}
\section{Artenerfassung} 
Die Arten eines Areals können auf unterschiedliche Arten und Weisen erfasst werden; betrachtet werden \emph{welche} Arten (qualitativ) und \emph{wie viele} Lebewesen dieser Arten (quantitativ) aufzufinden sind.

Dies kann mit Lupen oder Mikroskopen, einem Lupenbecher, dem bloßen Auge und so weiter geschehen, je nach größe des Lebewesens gemacht werden. So können die Lebewesens eines bestimmten Ortes betrachtet, klassifiziert, gezählt und bei Bedarf skizziert werden.

\subsection{Auswertung}
Anhand der erfassten Arten können Aussagen zur Dominanz mancher Arten gemacht werden; Lebewesen, von denen es mehr gibt, sind in der Regel dominanter, oder es können Aussagen zur Biodiversität gemacht werden; diese ist größer wenn es in einem Areal mehr unterschiedliche Arten gibt.

Zudem können generelle Aussagen zu den Umweltfaktoren gemacht werden, wenn alle gefundenen Lebewesen in Bezug auf einen bestimmten Umweltfaktor eine ähnliche Nische haben kann davon ausgegangen werden, dass dieser Umweltfaktor in dem Areal, zumindest langzeitig gesehen, vorliegt.
\end{document}